% !TeX spellcheck = es_CL
\documentclass[twocolumn]{article}
\usepackage[margin=1.2cm]{geometry}
\usepackage{tikz,lipsum,amsmath,amssymb,soul,xcolor,cancel,esint}
\pagestyle{empty}
\setlength{\columnseprule}{0.4pt}
\tikzset{every picture/.style={line width=0.75pt}}    

% opening
\title{\vspace{-4ex}\textbf{Formulario: Electrodinámica Avanzada}\vspace{-4ex}}
\author{\vspace{-5ex}}
\date{\vspace{-5ex}}
\pagenumbering{gobble}
\begin{document}
	\maketitle
	
	% ==========================================
	% BLOQUE 1: ONDAS PLANAS Y POLARIZACIÓN
	% ==========================================
	\begin{center}\textbf{\footnotesize ONDAS PLANAS Y POLARIZACIÓN}\end{center}
	\vspace{-2mm}
	\begin{equation*}
		\mathbf{E}(\mathbf{r},t) = \text{Re}\left[\mathbf{E}_0 e^{i(\mathbf{k}\cdot\mathbf{r} - \omega t)}\right], \quad \mathbf{B}(\mathbf{r},t) = \text{Re}\left[\mathbf{B}_0 e^{i(\mathbf{k}\cdot\mathbf{r} - \omega t)}\right]
	\end{equation*}
	\begin{equation*}
		\nabla^2\mathbf{E} - \mu\epsilon \frac{\partial^2\mathbf{E}}{\partial t^2} = 0, \quad \nabla^2\mathbf{H} - \mu\epsilon \frac{\partial^2\mathbf{H}}{\partial t^2} = 0
	\end{equation*}
	\begin{equation*}
		v = \frac{1}{\sqrt{\mu\epsilon}}, \quad n = \frac{c}{v} = \sqrt{\frac{\epsilon\mu}{\epsilon_0\mu_0}}, \quad Z = \sqrt{\frac{\mu}{\epsilon}} = \frac{Z_0}{n}
	\end{equation*}
	\begin{equation*}
		\mathbf{k}\cdot\mathbf{E}_0 = 0, \quad \mathbf{k}\cdot\mathbf{B}_0 = 0, \quad k^2 = \mu\epsilon\omega^2
	\end{equation*}
	\begin{equation*}
		\mathbf{B}_0 = \frac{1}{\omega}\mathbf{k}\times\mathbf{E}_0, \quad \mathbf{H}_0 = \frac{1}{Z}\hat{\mathbf{k}}\times\mathbf{E}_0
	\end{equation*}
	\begin{equation*}
		\langle u \rangle = \frac{1}{4}\epsilon |\mathbf{E}_0|^2 + \frac{1}{4\mu}|\mathbf{B}_0|^2 = \frac{1}{2}\epsilon |\mathbf{E}_0|^2
	\end{equation*}
	\begin{equation*}
		\mathbf{S} = \mathbf{E}\times\mathbf{H}, \quad \langle\mathbf{S}\rangle = \frac{1}{2}\text{Re}(\mathbf{E}_0\times\mathbf{H}_0^*) = \frac{|\mathbf{E}_0|^2}{2Z}\hat{\mathbf{k}}
	\end{equation*}
	\begin{equation*}
		\mathbf{J} = \begin{pmatrix} E_x \\ E_y \end{pmatrix} = \begin{pmatrix} a_x e^{i\delta_x} \\ a_y e^{i\delta_y} \end{pmatrix}, \quad \Delta = \delta_y - \delta_x
	\end{equation*}
	\begin{equation*}
		\frac{E_x^2}{a_x^2} + \frac{E_y^2}{a_y^2} - 2\frac{E_x E_y}{a_x a_y}\cos\Delta = \sin^2\Delta
	\end{equation*}
	\begin{align*}
		S_0 = a_x^2 + a_y^2, & \quad S_1 = a_x^2 - a_y^2 \\
		S_2 = 2a_x a_y\cos\Delta, & \quad S_3 = 2a_x a_y\sin\Delta
	\end{align*}
	\begin{equation*}
		S_0^2 \geq S_1^2 + S_2^2 + S_3^2 \quad (\text{Esfera de Poincaré})
	\end{equation*}
	
	\vspace{-2mm}\hrule\vspace{1mm}
	
	% ==========================================
	% BLOQUE 2: INTERFACES Y FRESNEL
	% ==========================================
	\begin{center}\textbf{\footnotesize INTERFACES Y COEFICIENTES DE FRESNEL}\end{center}
	\vspace{-2mm}
	\begin{equation*}
		\theta_r = \theta_i, \quad n_1\sin\theta_i = n_2\sin\theta_t \quad (\text{Snell})
	\end{equation*}
	\begin{equation*}
		r_s = \frac{n_1\cos\theta_i - n_2\cos\theta_t}{n_1\cos\theta_i + n_2\cos\theta_t}, \quad t_s = \frac{2n_1\cos\theta_i}{n_1\cos\theta_i + n_2\cos\theta_t}
	\end{equation*}
	\begin{equation*}
		r_p = \frac{n_2\cos\theta_i - n_1\cos\theta_t}{n_2\cos\theta_i + n_1\cos\theta_t}, \quad t_p = \frac{2n_1\cos\theta_i}{n_2\cos\theta_i + n_1\cos\theta_t}
	\end{equation*}
	\begin{equation*}
		R_{s,p} = |r_{s,p}|^2, \quad T_{s,p} = \frac{n_2\cos\theta_t}{n_1\cos\theta_i}|t_{s,p}|^2, \quad R+T=1
	\end{equation*}
	\begin{equation*}
		\tan\theta_B = \frac{n_2}{n_1} \implies r_p = 0 \quad (\text{Ángulo de Brewster})
	\end{equation*}
	\begin{equation*}
		\theta_i > \theta_c = \arcsin\left(\frac{n_2}{n_1}\right), \quad \kappa = \frac{\omega}{c}\sqrt{n_1^2\sin^2\theta_i - n_2^2}
	\end{equation*}
	
	\vspace{-2mm}\hrule\vspace{1mm}
	
	% ==========================================
	% BLOQUE 3: MEDIOS DISPERSIVOS Y PLASMAS
	% ==========================================
	\begin{center}\textbf{\footnotesize MEDIOS DISPERSIVOS, CONDUCTORES Y PLASMAS}\end{center}
	\vspace{-2mm}
	\begin{equation*}
		\nabla^2\mathbf{E} - \mu\epsilon\frac{\partial^2\mathbf{E}}{\partial t^2} - \mu\sigma\frac{\partial\mathbf{E}}{\partial t} = 0 \quad (\mathbf{J}_f = \sigma\mathbf{E})
	\end{equation*}
	\begin{equation*}
		\epsilon_c(\omega) = \epsilon + i\frac{\sigma}{\omega}, \quad k_c^2 = \mu\epsilon_c(\omega)\omega^2, \quad k_c = \beta + i\frac{\alpha}{2}
	\end{equation*}
	\begin{equation*}
		\sigma \gg \omega\epsilon \implies \beta \approx \frac{\alpha}{2} \approx \sqrt{\frac{\mu\sigma\omega}{2}}, \quad \delta = \frac{2}{\alpha} = \sqrt{\frac{2}{\mu\sigma\omega}}
	\end{equation*}
	\begin{equation*}
		\omega_p = \sqrt{\frac{N e^2}{m \epsilon_0}}, \quad \epsilon(\omega) = \epsilon_0\left(1 - \frac{\omega_p^2}{\omega^2}\right), \quad \omega^2 = \omega_p^2 + c^2 k^2
	\end{equation*}
	\begin{equation*}
		v_\phi = \frac{\omega}{k} = \frac{c}{\sqrt{1 - \omega_p^2/\omega^2}}, \quad v_g = \frac{d\omega}{dk} = c\sqrt{1 - \frac{\omega_p^2}{\omega^2}}
	\end{equation*}
	\begin{equation*}
		v_\phi v_g = c^2, \quad k(\omega) \approx k(\omega_0) + \frac{1}{v_g}(\omega - \omega_0) + \frac{1}{2}k''(\omega_0)(\omega - \omega_0)^2
	\end{equation*}
	
	\vspace{-2mm}\hrule\vspace{1mm}
	
	% ==========================================
	% BLOQUE 4: CAUSALIDAD Y KRAMERS-KRONIG
	% ==========================================
	\begin{center}\textbf{\footnotesize CAUSALIDAD Y KRAMERS-KRONIG}\end{center}
	\vspace{-2mm}
	\begin{equation*}
		\mathbf{P}(t) = \epsilon_0 \int_{-\infty}^{t} \chi(t - t')\mathbf{E}(t') dt' \implies \chi(-\omega) = \chi^*(\omega)
	\end{equation*}
	\begin{equation*}
		\text{Re}\,\chi(\omega) = \frac{1}{\pi}\mathcal{P}\int_{-\infty}^{\infty} \frac{\text{Im}\,\chi(\omega')}{\omega' - \omega} d\omega'
	\end{equation*}
	\begin{equation*}
		\text{Im}\,\chi(\omega) = -\frac{1}{\pi}\mathcal{P}\int_{-\infty}^{\infty} \frac{\text{Re}\,\chi(\omega')}{\omega' - \omega} d\omega'
	\end{equation*}
	\begin{equation*}
		\text{Re}\,\epsilon(\omega) - 1 = \frac{2}{\pi}\mathcal{P}\int_{0}^{\infty} \frac{\omega' \text{Im}\,\epsilon(\omega')}{\omega'^2 - \omega^2} d\omega'
	\end{equation*}
	\begin{equation*}
		\text{Im}\,\epsilon(\omega) = -\frac{2\omega}{\pi}\mathcal{P}\int_{0}^{\infty} \frac{\text{Re}\,\epsilon(\omega') - 1}{\omega'^2 - \omega^2} d\omega'
	\end{equation*}
	\begin{equation*}
		\epsilon(\omega) = 1 + \int_{0}^{\infty} \frac{f(\Omega)}{\Omega^2 - \omega^2 - i\Gamma\omega} d\Omega
	\end{equation*}
	\begin{equation*}
		\lim_{\Gamma \to 0^+} \text{Im}\,\epsilon(\omega) = \frac{\pi}{2\omega} f(\omega) \quad (\omega > 0)
	\end{equation*}
	
	\vspace{-2mm}\hrule\vspace{1mm}
	
	% ==========================================
	% BLOQUE 5: GUÍAS DE ONDA
	% ==========================================
	\begin{center}\textbf{\footnotesize GUÍAS DE ONDA}\end{center}
	\vspace{-2mm}
	\begin{equation*}
		\mathbf{E} = \left[\mathbf{E}_\perp + E_z\hat{\mathbf{z}}\right]e^{i(kz-\omega t)}, \quad \mathbf{B} = \left[\mathbf{B}_\perp + B_z\hat{\mathbf{z}}\right]e^{i(kz-\omega t)}
	\end{equation*}
	\begin{equation*}
		(\nabla_\perp^2 + \gamma^2)E_z = 0, \quad (\nabla_\perp^2 + \gamma^2)B_z = 0, \quad \gamma^2 = \frac{\omega^2}{c^2} - k^2
	\end{equation*}
	\begin{equation*}
		\mathbf{E}_\perp = \frac{i}{\gamma^2}\left(k\nabla_\perp E_z - \omega\hat{\mathbf{z}}\times\nabla_\perp B_z\right)
	\end{equation*}
	\begin{equation*}
		\mathbf{B}_\perp = \frac{i}{\gamma^2}\left(k\nabla_\perp B_z + \frac{\omega}{c^2}\hat{\mathbf{z}}\times\nabla_\perp E_z\right)
	\end{equation*}
	\begin{equation*}
		\text{TM: } B_z=0, \ \left.E_z\right|_S = 0; \quad \text{TE: } E_z=0, \ \left.\frac{\partial B_z}{\partial n}\right|_S = 0
	\end{equation*}
	\begin{equation*}
		\text{Guía Rec: } \omega_{mn} = c\gamma_{mn} = c\sqrt{\left(\frac{m\pi}{a}\right)^2 + \left(\frac{n\pi}{b}\right)^2}
	\end{equation*}
	\begin{equation*}
		k_{mn}(\omega) = \frac{1}{c}\sqrt{\omega^2 - \omega_{mn}^2}, \quad v_\phi v_g = c^2
	\end{equation*}
	\begin{equation*}
		Z_{\text{TE}} = \frac{Z_0}{\sqrt{1 - (\omega_{mn}/\omega)^2}}, \quad Z_{\text{TM}} = Z_0\sqrt{1 - \left(\frac{\omega_{mn}}{\omega}\right)^2}
	\end{equation*}
	\begin{equation*}
		\alpha = \frac{P_{\text{dis}}}{2 P_{\text{trans}}}, \quad P_{\text{dis}} = \frac{1}{2\sigma\delta}\oint_{S} |\mathbf{H}_{\text{tang}}|^2 dl
	\end{equation*}
	
	\vspace{-2mm}\hrule\vspace{1mm}
	
	% ==========================================
	% BLOQUE 6: RADIACIÓN MULTIPOLAR
	% ==========================================
	\begin{center}\textbf{\footnotesize POTENCIALES RETARDADOS Y RADIACIÓN}\end{center}
	\vspace{-2mm}
	\begin{equation*}
		\mathbf{A}(\mathbf{r}) = \frac{\mu_0}{4\pi}\int \mathbf{J}(\mathbf{r}')\frac{e^{ik|\mathbf{r}-\mathbf{r}'|}}{|\mathbf{r}-\mathbf{r}'|} d^3r'
	\end{equation*}
	\begin{equation*}
		r \gg \lambda \implies \mathbf{A}(\mathbf{r}) \approx \frac{\mu_0}{4\pi}\frac{e^{ikr}}{r}\int \mathbf{J}(\mathbf{r}')e^{-ik\hat{\mathbf{r}}\cdot\mathbf{r}'}d^3r'
	\end{equation*}
	\begin{equation*}
		\mathbf{B} \approx ik\hat{\mathbf{r}}\times\mathbf{A}, \quad \mathbf{E} \approx c\mathbf{B}\times\hat{\mathbf{r}}
	\end{equation*}
	\begin{equation*}
		\mathbf{p} = \int \mathbf{r}'\rho(\mathbf{r}')d^3r', \quad \mathbf{A}(\mathbf{r}) = - \frac{i\mu_0\omega}{4\pi}\mathbf{p}\frac{e^{ikr}}{r}
	\end{equation*}
	\begin{equation*}
		\mathbf{B} = \frac{\mu_0\omega^2}{4\pi c}(\hat{\mathbf{r}}\times\mathbf{p})\frac{e^{ikr}}{r}, \quad \mathbf{E} = \frac{\mu_0\omega^2}{4\pi}\left[(\hat{\mathbf{r}}\times\mathbf{p})\times\hat{\mathbf{r}}\right]\frac{e^{ikr}}{r}
	\end{equation*}
	\begin{equation*}
		\frac{dP}{d\Omega} = \frac{\mu_0\omega^4}{32\pi^2 c}|\mathbf{p}|^2\sin^2\theta, \quad P_{\text{total}} = \frac{\mu_0\omega^4}{12\pi c}|\mathbf{p}|^2
	\end{equation*}
	\begin{equation*}
		R_{\text{rad}} = \frac{2 P_{\text{total}}}{I_0^2} = \frac{\mu_0\omega^2 d^2}{6\pi c} \quad (\text{Dipolo eléctrico } E1)
	\end{equation*}
	
	\vspace{-2mm}\hrule\vspace{1mm}
	
	% ==========================================
	% BLOQUE 7: ELECTRODINÁMICA COVARIANTE
	% ==========================================
	\begin{center}\textbf{\footnotesize ELECTRODINÁMICA COVARIANTE}\end{center}
	\vspace{-2mm}
	\begin{equation*}
		\eta_{\mu\nu} = \eta^{\mu\nu} = \text{diag}(1, -1, -1, -1)
	\end{equation*}
	\begin{equation*}
		x^\mu = (ct, \mathbf{r}), \quad \partial_\mu = \left(\frac{1}{c}\frac{\partial}{\partial t}, \nabla\right), \quad \square = \partial_\mu\partial^\mu = \frac{1}{c^2}\frac{\partial^2}{\partial t^2} - \nabla^2
	\end{equation*}
	\begin{equation*}
		U^\mu = \gamma(c, \mathbf{v}), \quad p^\mu = m U^\mu = (E/c, \mathbf{p}), \quad p_\mu p^\mu = m^2 c^2
	\end{equation*}
	\begin{equation*}
		J^\mu = (c\rho, \mathbf{J}), \quad A^\mu = (\Phi/c, \mathbf{A}), \quad F_{\mu\nu} = \partial_\mu A_\nu - \partial_\nu A_\mu
	\end{equation*}
	\begin{equation*}
		F^{\mu\nu} = \begin{pmatrix} 0 & -E_x/c & -E_y/c & -E_z/c \\ E_x/c & 0 & -B_z & B_y \\ E_y/c & B_z & 0 & -B_x \\ E_z/c & -B_y & B_x & 0 \end{pmatrix}
	\end{equation*}
	\begin{equation*}
		\tilde{F}^{\mu\nu} = \frac{1}{2}\epsilon^{\mu\nu\rho\sigma}F_{\rho\sigma}, \quad (\mathbf{E}/c \to \mathbf{B}, \ \mathbf{B} \to -\mathbf{E}/c)
	\end{equation*}
	\begin{equation*}
		\partial_\mu F^{\mu\nu} = \mu_0 J^\nu, \quad \partial_\mu \tilde{F}^{\mu\nu} = 0, \quad \partial_\mu J^\mu = 0
	\end{equation*}
	\begin{equation*}
		\partial_\mu A^\mu = 0 \implies \square A^\mu = \mu_0 J^\mu \quad (\text{Gauge de Lorenz})
	\end{equation*}
	\begin{equation*}
		F_{\mu\nu}F^{\mu\nu} = 2\left(B^2 - \frac{E^2}{c^2}\right), \quad F_{\mu\nu}\tilde{F}^{\mu\nu} = -\frac{4}{c}(\mathbf{E}\cdot\mathbf{B})
	\end{equation*}
	
	\vspace{-2mm}\hrule\vspace{1mm}
	
	% ==========================================
	% BLOQUE 8: DINÁMICA RELATIVISTA Y DRIFTS
	% ==========================================
	\begin{center}\textbf{\footnotesize DINÁMICA RELATIVISTA Y DRIFTS}\end{center}
	\vspace{-2mm}
	\begin{align*}
		E'_x = E_x, & \quad B'_x = B_x \\
		E'_y = \gamma(E_y - v B_z), & \quad B'_y = \gamma\left(B_y + \frac{v}{c^2}E_z\right) \\
		E'_z = \gamma(E_z + v B_y), & \quad B'_z = \gamma\left(B_z - \frac{v}{c^2}E_y\right)
	\end{align*}
	\begin{equation*}
		K^\mu = \frac{dp^\mu}{d\tau} = q F^{\mu\nu}U_\nu \implies \frac{d\mathbf{p}}{dt} = q(\mathbf{E} + \mathbf{v}\times\mathbf{B})
	\end{equation*}
	\begin{equation*}
		L = -mc^2\sqrt{1 - \frac{v^2}{c^2}} - q\Phi + q\mathbf{A}\cdot\mathbf{v}
	\end{equation*}
	\begin{equation*}
		\mathbf{P}_{\text{can}} = \mathbf{p} + q\mathbf{A}, \quad H = \sqrt{(\mathbf{P} - q\mathbf{A})^2 c^2 + m^2 c^4} + q\Phi
	\end{equation*}
	\begin{equation*}
		\mathbf{v}_E = \frac{\mathbf{E}\times\mathbf{B}}{B^2}, \quad \mu = \frac{m v_\perp^2}{2B} = \text{const} \implies \frac{B_{\text{esp}}}{B_0} = \frac{v_0^2}{v_{\perp 0}^2}
	\end{equation*}
	
	\vspace{-2mm}\hrule\vspace{1mm}
	
	% ==========================================
	% BLOQUE 9: ACCIÓN Y TENSORES DE ENERGÍA
	% ==========================================
	\begin{center}\textbf{\footnotesize ACCIÓN Y TENSORES DE ENERGÍA}\end{center}
	\vspace{-2mm}
	\begin{equation*}
		\mathcal{L}_{\text{Maxwell}} = -\frac{1}{4\mu_0}F_{\mu\nu}F^{\mu\nu} - J_\mu A^\mu
	\end{equation*}
	\begin{equation*}
		\mathcal{L}_{\text{Proca}} = -\frac{1}{4\mu_0}F_{\mu\nu}F^{\mu\nu} + \frac{\mu_\gamma^2}{2\mu_0}A_\mu A^\mu \quad (\text{Masa fotón})
	\end{equation*}
	\begin{equation*}
		T_{\text{can}}^{\mu\nu} = -\frac{1}{4\mu_0}\eta^{\mu\nu}F_{\alpha\beta}F^{\alpha\beta} - \frac{1}{\mu_0}F^{\mu\rho}\partial^\nu A_\rho
	\end{equation*}
	\begin{equation*}
		\Theta^{\mu\nu} = T_{\text{can}}^{\mu\nu} + \frac{1}{\mu_0}F^{\mu\rho}\partial_\rho A^\nu = \frac{1}{\mu_0}\left(F^{\mu\rho}F^\nu{}_\rho + \frac{1}{4}\eta^{\mu\nu}F_{\alpha\beta}F^{\alpha\beta}\right)
	\end{equation*}
	
	\subsection*{1. Guías de Onda Rectangulares ($0 \le x \le a$, $0 \le y \le b$)}
	Ecuación de Helmholtz: $(\nabla_\perp^2 + \gamma^2) \psi = 0$, con $\gamma^2 = \frac{\omega^2}{c^2} - k^2$.
	\begin{enumerate}
		\item {\bf Modo TM:} $B_z = 0$, hallar $E_z$. Condición de Dirichlet: $E_z|_{\partial\Omega} = 0$.
		$$E_z(x,y) = E_0 \sin\left(\frac{m\pi x}{a}\right)\sin\left(\frac{n\pi y}{b}\right) \quad (m,n \ge 1)$$
		\item {\bf Modo TE:} $E_z = 0$, hallar $B_z$. Condición de Neumann: $\frac{\partial B_z}{\partial n}|_{\partial\Omega} = 0$.
		$$B_z(x,y) = B_0 \cos\left(\frac{m\pi x}{a}\right)\cos\left(\frac{n\pi y}{b}\right) \quad (m,n \ge 0, m+n > 0)$$
		\item {\bf Corte:} Frecuencia de corte límite $\omega_{mn} = c \sqrt{(\frac{m\pi}{a})^2 + (\frac{n\pi}{b})^2}$.
		\item {\bf Transversales:} Hallar $\mathbf{E}_\perp, \mathbf{B}_\perp$ derivando componentes longitudinales:
		$$\mathbf{E}_\perp = \frac{1}{\gamma^2} [ik\nabla_\perp E_z - i\omega(\hat{z}\times\nabla_\perp B_z)]$$
	\end{enumerate}
	
	\subsection*{2. Transformación de Campos ($F^{\mu\nu}$)}
	Métrica $\eta_{\mu\nu}=\text{diag}(1,-1,-1,-1)$. Transformación tensorial: $F' = \Lambda F \Lambda^T$.
	\begin{enumerate}
		\item {\bf Tensor $F^{\mu\nu}$:} $F^{0i} = -E_i/c$, $F^{ij} = -\epsilon_{ijk}B_k$.
		\item {\bf Boost (eje $x$):} $\Lambda^0_{\ 0}=\gamma, \Lambda^0_{\ 1}=\Lambda^1_{\ 0}=-\beta\gamma, \Lambda^1_{\ 1}=\gamma$.
		\item {\bf Campos Rotados / Transformados:}
		$$E'_x = E_x, \quad E'_y = \gamma(E_y - vB_z), \quad E'_z = \gamma(E_z + vB_y)$$
		$$B'_x = B_x, \quad B'_y = \gamma(B_y + \frac{v}{c^2}E_z), \quad B'_z = \gamma(B_z - \frac{v}{c^2}E_y)$$
		\item {\bf Invariantes:} Evaluar $E^2-c^2B^2$ y $\mathbf{E}\cdot\mathbf{B}$ para chequear consistencia.
	\end{enumerate}
	
	\subsection*{3. Radiación de Dipolo Eléctrico Oscilante}
	Fuente armónica $\rho(\mathbf{r},t)=\rho_0(\mathbf{r})e^{-i\omega t}$. Zona lejana ($r \gg \lambda \gg d$).
	\begin{enumerate}
		\item {\bf Momento Dipolar:} $\mathbf{p} = \int \mathbf{r}\rho_0(\mathbf{r})d^3r$.
		\item {\bf Campos de Radiación:} $\mathbf{B} = \frac{\mu_0\omega^2}{4\pi c}(\hat{r}\times\mathbf{p})\frac{e^{ikr}}{r}$ y $\mathbf{E} = c(\mathbf{B}\times\hat{r})$.
		\item {\bf Vector de Poynting:} $\langle\mathbf{S}\rangle = \frac{\mu_0\omega^4|p|^2}{32\pi^2 c}\frac{\sin^2\theta}{r^2}\hat{r}$ (si $\mathbf{p}=p\hat{z}$).
		\item {\bf Potencia Total:} Integral de flujo angular $\implies P = \frac{\mu_0\omega^4|p|^2}{12\pi c}$.
	\end{enumerate}
	
	\subsection*{4. Relaciones de Kramers-Kronig}
	Consecuencia directa de la causalidad lineal ($\chi(t)=0$ para $t<0$).
	$$\text{Re }\epsilon(\omega) - 1 = \frac{2}{\pi}\mathcal{P}\int_0^\infty \frac{\Omega \text{Im }\epsilon(\Omega)}{\Omega^2-\omega^2}d\Omega$$
	
\end{document}